% year: תשפ"ב
% type: מבחן
% semester: א
% moed: מיוחד
% number: 5
% points: 17
% categories: antiderivatives
% source: exams/מבחנים/תשפב/מועד ג תשפב.pdf

\begin{question}
חשבו את האינטגרל
\[ \int\frac{5x^2+2x+2}{x^3-1}dx \]
\end{question}

\begin{hint}
$x^3-1=(x-1)(x^2+x+1)$, והגורם הריבועי אי-פריק. פרקו לשברים חלקיים.
\end{hint}

\begin{solution}
\textbf{פירוק לשברים חלקיים.} $x^3-1=(x-1)(x^2+x+1)$, ולגורם $x^2+x+1$ דיסקרימיננטה $1-4<0$, ולכן הוא אי-פריק. מעלת המונה קטנה ממעלת המכנה:
\[ \frac{5x^2+2x+2}{(x-1)(x^2+x+1)}=\frac{A}{x-1}+\frac{Bx+C}{x^2+x+1}, \]
\[ 5x^2+2x+2=A(x^2+x+1)+(Bx+C)(x-1). \]
\begin{itemize}
  \item הצבת $x=1$: $9=3A$, ולכן $A=3$.
  \item אז $5x^2+2x+2-3(x^2+x+1)=2x^2-x-1=(x-1)(2x+1)$, ולכן $Bx+C=2x+1$, כלומר $B=2$, $C=1$.
\end{itemize}
\[ \frac{5x^2+2x+2}{x^3-1}=\frac{3}{x-1}+\frac{2x+1}{x^2+x+1}. \]

\textbf{אינטגרציה.} $\int\frac{3}{x-1}dx=3\ln|x-1|$, והמונה $2x+1$ הוא בדיוק נגזרת המכנה $x^2+x+1>0$, ולכן $\int\frac{2x+1}{x^2+x+1}dx=\ln(x^2+x+1)$.

\textbf{תשובה:}
\[ \int\frac{5x^2+2x+2}{x^3-1}dx=\boxed{3\ln|x-1|+\ln(x^2+x+1)+C}. \]
\end{solution}
